\documentclass[12pt]{article}
\usepackage[utf8]{inputenc}
\usepackage[T1]{fontenc}
\usepackage[spanish]{babel}
\usepackage{amsmath, amssymb, amsfonts}
\usepackage{fancyhdr}
\usepackage{xcolor}
\usepackage[margin=3cm]{geometry}
\usepackage{graphicx}
\usepackage{float}
\usepackage{multicol}
\usepackage{enumitem}
\usepackage{tcolorbox}
\usepackage{listings}

% Configuración de listings para Python
\definecolor{codegreen}{rgb}{0,0.6,0}
\definecolor{codegray}{rgb}{0.5,0.5,0.5}
\definecolor{codepurple}{rgb}{0.58,0,0.82}
\definecolor{backcolour}{rgb}{0.95,0.95,0.92}

\lstdefinestyle{mystyle}{
    backgroundcolor=\color{backcolour},   
    commentstyle=\color{codegreen},
    keywordstyle=\color{magenta},
    numberstyle=\tiny\color{codegray},
    stringstyle=\color{codepurple},
    basicstyle=\ttfamily\small,
    breakatwhitespace=false,         
    breaklines=true,                 
    captionpos=b,                    
    keepspaces=true,                 
    numbers=left,                    
    numbersep=5pt,                  
    showspaces=false,                
    showstringspaces=false,
    showtabs=false,                  
    tabsize=2,
    language=Python
}

\lstset{style=mystyle}

% Encabezado
\pagestyle{fancy}
\fancyhf{}
\fancyhead[L]{Escuela de Ingeniería \\ Universidad de O'Higgins}
\fancyhead[R]{Programación\\ 2026 01}
\fancyfoot[C]{\thepage}

% Título comprimido
\makeatletter
\renewcommand{\maketitle}{%
  \begin{center}
    {\LARGE \@title\par}
    \vskip 0.3em
    {\large \@author\par}
    \vskip 0.3em
  \end{center}
  \vskip -0.2em  % Espacio después del título
}
\makeatother

\title{Guía repaso de funciones}
\author{Felipe Llach R.}

\begin{document}
\maketitle

\section{Introducción}
Las funciones son bloques de código reutilizables que realizan una tarea específica. Nos permiten organizar el código, evitar la repetición y facilitar la resolución de problemas complejos dividiéndolos en partes más pequeñas.

\section{Definición de una Función}
Para definir una función en Python se utiliza la palabra reservada \texttt{def}, seguida del nombre de la función y paréntesis.

\begin{tcolorbox}[title=Estructura básica]
\begin{lstlisting}
def nombre_de_la_funcion(parametros):
    # Cuerpo de la funcion
    instrucciones
    return resultado # Opcional
\end{lstlisting}
\end{tcolorbox}

\section{Parámetros y Argumentos}
Los \textbf{parámetros} son las variables que se listan en la definición. Los \textbf{argumentos} son los valores reales que se pasan al llamar a la función.

\subsection{Valores por Defecto}
Es posible asignar valores por defecto a los parámetros. Si al llamar a la función no se proporciona un argumento para ese parámetro, se utilizará el valor predefinido.

\begin{lstlisting}
def saludar(nombre, mensaje="Hola"):
    print(f"{mensaje}, {nombre}!")

saludar("Diego")           # Imprime: Hola, Diego!
saludar("Elena", "Buenos dias") # Imprime: Buenos dias, Elena!
\end{lstlisting}

\section{Retorno de Valores}
La sentencia \texttt{return} permite que la función envíe un resultado de vuelta al lugar donde fue llamada.

\subsection{Retorno de múltiples valores}
Python permite retornar más de un valor a la vez. Técnicamente, los valores se retornan como una \textbf{tupla}, que luego puede ser "desempaquetada".

\begin{lstlisting}
def calcular_rectangulo(base, altura):
    area = base * altura
    perimetro = 2 * (base + altura)
    return area, perimetro

# Desempaquetando los valores
a, p = calcular_rectangulo(10, 5)
print(f"Area: {a}, Perimetro: {p}")
\end{lstlisting}

\section{Documentación de Funciones (Docstrings)}
Se utilizan triples comillas (\texttt{"""}) para documentar la función. Se puede consultar con \texttt{help()}.

\begin{lstlisting}
def sumar(a, b):
    """Retorna la suma de a y b."""
    return a + b

help(sumar)
\end{lstlisting}

\section{Módulos e Importación}
Para usar funciones de otros archivos (\texttt{.py}) o librerías estándar:

\begin{lstlisting}
import math
from random import randint

print(math.sqrt(25))
print(randint(1, 100))
\end{lstlisting}

\end{document}
